\documentclass[10pt,a4paper]{amsart}
\usepackage[margin=19mm]{geometry}
\usepackage{amsmath,amssymb,mathtools}
\usepackage[scr=rsfs,cal=euler]{mathalfa}
\usepackage{xfrac}
\usepackage[unicode,hidelinks,bookmarks=false]{hyperref}
\newcommand{\ie}{i.e.\ }
\newcommand{\cf}{cf.\ }
\newcommand{\resp}{respectively\ }
\newcommand{\Subsection}{\subsection}
\providecommand{\nobreakdash}{\nobreak\hbox{-}}
\setlength{\emergencystretch}{2em}
\multlinegap0pt
\usepackage{bm}
\usepackage{amssymb}
\usepackage{mathtools}
\usepackage{algorithm}\usepackage{algpseudocode}
\newtheorem{theorem}{Theorem}
\title{Message-Passing GNNs Fail to Approximate \\ Sparse Triangular Factorizations}
\author{Vladislav Trifonov, Ekaterina Muravleva, Ivan Oseledets}
\date{Source: arXiv:2502.01397v3 (CC BY 4.0)}
\begin{document}
\maketitle

\noindent\emph{Source and licence.} Message-Passing GNNs Fail to Approximate \\ Sparse Triangular Factorizations. Version arXiv:2502.01397v3 (CC BY 4.0), distributed by arXiv under the Creative Commons Attribution 4.0 International licence, http://creativecommons.org/licenses/by/4.0/, as recorded in the arXiv licence metadata for this exact version and shown on the abstract page https://arxiv.org/abs/2502.01397v3. Full licence notice: \texttt{licenses/arxiv-2502.01397-CC-BY-4.0.txt}.\par\smallskip

\noindent\textit{Editorial scope.} Counted unit: the article's theorem theorem (source0.tex lines 662--670). The article's complete original proof of it is delivered with it (source0.tex lines 673--686, one \texttt{proof} environment of 14 lines). No mathematics is written, repaired or re-derived: the statement and the proof are the article's own lines, in the article's own notation, and no retained line is substituted or re-flowed.  No further source-local context is needed: every cross-reference of every retained line resolves inside the two delivered spans.  Nothing was omitted from any retained span, so the package carries no omission record.  No retained line contains a citation, so the wrapper carries no bibliography and no bibliography entry.  No retained line contains a figure environment, an image-inclusion command or drawing code, so no image is delivered and the package declares no asset dependency.  Editorial formatting only, in addition: an AMS \texttt{amsart} preamble that restates the article's own declarations and macros used by the retained lines, namely the environments \texttt{theorem} on separate counters, together with the standard packages those lines need.  The retained environments print in this document as Theorem 1 in order of appearance, whereas the article prints the same items with the numbering of its own full document.  The complete native source of the article as distributed in the arXiv e-print for this exact version is bundled unchanged as \texttt{sources/172653/source0.tex}.
\begin{theorem}
    Let $A$ be a tridiagonal symmetric positive definite $n \times n$ matrix. Then it can be factorized as
    \begin{equation*}
        A = LL^{\top},
    \end{equation*}
    where $L$ is a bidiagonal lower triangular matrix, and then mapping $A \rightarrow L$ is not local, which means 
    that there exist matrix $A$ and $A'$ such that $A-A'$ has only one non-zero element, where
    the difference $L - L'$ has dense support: many entries change significantly, even though only a single entry was perturbed.
\end{theorem}

\begin{proof}
    Consider the matrix $A$ given by $A=LL^{\top}$ where $L$ is a bidiagonal matrix with $L_{ii} = \frac{1}{i}, i = 1, \ldots, n$ and
    $L_{i, i-1} = 1, i = 2$. Then $A$ is a symmetric positive definite tridiagonal matrix with elements $A_{11} = 1, A_{i, i} = 1 + \frac{1}{i^2}, 
    A_{i+1, i} = A_{i, i+1} = \frac{1}{i}, i = 1, \ldots, n-1$.
    Now, consider the matrix $A' = A + e_1 e_1^{\top}$, where $e_1$ is the first column of the identity matrix. Let
    $A' = L' L'^{\top}$ be its Cholesky factorization. The matrix $L'$ is bidiagonal. The element $L'_{11}$ is equal to $\sqrt{2}$,
    and for each $i=2, \ldots, n$ we have the well-known formulas 
    \begin{equation*}
      L'_{i, i-1} = \frac{L_{i, i-1}}{L'_{i-1, i-1}} = \frac{\frac{1}{i-1}}{L'_{i-1, i-1}},
    \end{equation*} 
    and $L'_{i, i} = \sqrt{A_{i, i} - \left( L_{i, i-1}\right)^2 }.$  Let $d_i = (L'_{i, i})^2$, then
    $d_1 = 2, d_i = 1 + \frac{1}{i^2} - \frac{1}{d_{i-1} (i-1)^2}$. From this recurrence relation it is easy to see that $d_i$
    converges to $1$ as $i \to \infty$.
\end{proof}

\end{document}
